% Part: lambda-calculus
% Chapter: church-rosser
% Section: beta-eta-reduction

\documentclass[../../../include/open-logic-section]{subfiles}

\begin{document}

\olfileid{lam}{cr}{be}

\olsection{$\beta\eta$-راکمول}

د چرچ--روسر خاصيت د $\beta\eta$-راکمول ($\bered$) دپاره
هم صدق کوي۔ \textbf{د سرچينې سپيناوی (OLLAM-043):} په
لاندې ثبوت کښې د ګډ يوګامي راکمول نښه کارېږي، خو مخکې
يې اړيکه په څرګنده نۀ ده تعريف شوې؛ دلته مطلب د بېټا
يا اېټا د يوګامي انقباض اړونده اړيکه ده۔

\begin{lem}\ollabel{lem:one-par}
  که $M \beredone M'$ وي، نو $M \beredpar M'$.
\end{lem}

\begin{proof} 
  د $M \beredone M'$ د !!{derivation} پر جوړښت استقرا کوو۔
  که $M \eredone M'$ د $\eta$-بدلون له لارې وي، يعنې
  \olref[syn][eta]{defn:beredone}، نو
  \olref[pbe]{thm:refl} کاروو۔ نور حالتونه د
  \olref[b]{lem:one-par} په شان دي۔
  \textbf{د سرچينې سمون (OLLAM-044):} اصلي متن دلته
  د اېټا-بدلون دپاره د بېټا-انقباض نښه کاروي؛ د اړوندې
  اېټا نښې ته سمه شوې ده۔ \textbf{د ثبوت سپيناوی
  (OLLAM-045):} يوازې د سر رېډکس حالت بس نۀ دے؛
  د سازګاره اړيکې دننني حالتونه د ترم په جوړښت او
  د هممهاله قواعدو په کارولو پوره کېږي۔
\end{proof}


\begin{lem}\ollabel{lem:par-red}
  که $M \beredpar M'$ وي، نو $M \bered M'$.
\end{lem}

\begin{proof} د $M \beredpar M'$ د !!{derivation} پر جوړښت استقرا کوو۔

  که وروستۍ قاعده \olref[pbe]{defn:beredpar5} وي، نو $M$،
  $\lambd[x][Nx]$ او $M'$ د ځينو $x$، $N$ او $N'$ دپاره
  $N'$ دے، چې $x \notin FV(N)$ او $N \beredpar N'$ وي۔
  نو لومړے $\lambd[x][Nx]$ د $\eta$-بدلون له لارې $N$
  ته راکمولے شو؛ بيا د $\beredone$ د ګامونو هغه لړۍ
  تعقيبوو چې $N \bered N'$ ښيي۔ وروستۍ خبره د استقرا
  له فرضيې راځي۔
\end{proof}


\begin{lem}\ollabel{lem:str}
  $\bered$ تر ټولو کوچنۍ انتقالي اړيکه ده چې $\beredpar$
  پکې شامله ده۔
\end{lem}

\begin{proof}
  لکه څنګه چې په \olref[b]{lem:str} کښې وشول۔
\end{proof}

\begin{thm}\ollabel{thm:cr}
  $\bered$ د چرچ--روسر خاصيت لري۔
\end{thm}

\begin{proof}
  دا د \olref[dap]{thm:str}، \olref[pbe]{thm:cr} او
  \olref{lem:str} له مخې راځي۔
\end{proof}
\end{document}
